El sistema implementa tres algoritmos de planificación de CPU seleccionables de forma
interactiva al inicio de cada ejecución. Todos operan sobre la misma unidad de trabajo:
un archivo PDF de horario generado por el portal de Autoservicios BUAP. La llegada
de procesos equivale al momento en que el \texttt{DirectoryStream} entrega las rutas
al programa; a partir de ese instante, cada algoritmo decide el orden y la forma en
que la CPU atiende cada archivo.

La selección se realiza mediante un menú desplegado en consola antes del contador
de inicio, implementado en el método \texttt{elegirAlgoritmo()} de la clase
\texttt{Generador}:

\begin{figure}[H]
    \centering
    \caption{Menú interactivo de selección de algoritmo}
    \label{code:menu-algoritmo}
    \begin{lstlisting}[language=Java]
private AlgoritmoScheduler elegirAlgoritmo() {
    System.out.println("Algoritmo de planificacion de CPU");
    System.out.println("  [1] FCFS  - First Come, First Served");
    System.out.println("  [2] SJF   - Shortest Job First");
    System.out.println("  [3] RR    - Round Robin (quantum=150ms)");
    System.out.print("Selecciona una opcion [1-3]: ");

    Scanner scanner = new Scanner(System.in);
    String opcion = scanner.nextLine().trim();

    return switch (opcion) {
        case "1" -> new FCFSScheduler();
        case "2" -> new SJFScheduler();
        case "3" -> new RoundRobinScheduler();
        default  -> null;
    };
}
    \end{lstlisting}
\end{figure}

Los tres schedulers implementan la interfaz \texttt{AlgoritmoScheduler}, que define
dos métodos: \texttt{ordenar()}, que recibe la lista de PDFs en orden de llegada y
devuelve la lista reordenada según el criterio del algoritmo; y \texttt{ejecutar()},
que procesa los archivos y registra las métricas de cada uno.

% ─────────────────────────────────────────────────────────────────
\subsubsection{FCFS (First Come, First Served)}

FCFS es un algoritmo no expropiativo que atiende los procesos estrictamente en el
orden de llegada. En este sistema, la "llegada" equivale al orden en que el Sistema
Operativo entrega los archivos PDF a través de \texttt{DirectoryStream}. Una vez
que un archivo toma la CPU, la ocupa hasta finalizar su procesamiento completo.

La implementación no aplica ningún reordenamiento: el método \texttt{ordenar()}
devuelve la lista tal como fue recibida.

\begin{figure}[H]
    \centering
    \caption{Implementación de FCFS — sin reordenamiento}
    \label{code:fcfs-impl}
    \begin{lstlisting}[language=Java]
// FCFSScheduler.java
@Override
public List<Path> ordenar(List<Path> pdfs) {
    // No reordena: devuelve la lista tal como llego del SO
    return new ArrayList<>(pdfs);
}

@Override
public List<MetricaPDF> ejecutar(List<Path> pdfs) throws Exception {
    for (Path pdf : pdfs) {
        MetricaPDF metrica = new MetricaPDF(
            pdf.getFileName().toString(), tiempoInicioCola);
        metrica.registrarInicio();

        // El hilo ocupa la CPU hasta terminar (Non-preemptive)
        ProcesadorPDFTask tarea = new ProcesadorPDFTask(pdf);
        tarea.call();

        metrica.registrarFin();
        metricas.add(metrica);
    }
    return metricas;
}
    \end{lstlisting}
\end{figure}

El uso nativo de FCFS se presenta por la naturaleza secuencial de \texttt{DirectoryStream}.
El Sistema Operativo actúa como despachador, entregando las rutas en el orden en que
son recuperadas de la tabla de descriptores del directorio. Al no existir ninguna
evaluación heurística previa ni un temporizador que interrumpa la ejecución, el
tiempo de espera de cada archivo depende exclusivamente de su posición en la cola.

% ─────────────────────────────────────────────────────────────────
\subsubsection{SJF (Shortest Job First)}

SJF es un algoritmo no expropiativo que prioriza los procesos con menor burst time
estimado. En este sistema, el tamaño del archivo en bytes se utiliza como estimador:
un PDF más pequeño contiene menos contenido y se asume que requiere menor tiempo
de procesamiento por parte de la CPU.

Antes de ejecutar, el método \texttt{ordenar()} inspecciona el tamaño de cada
archivo con \texttt{Files.size()} y reordena la lista de menor a mayor:

\begin{figure}[H]
    \centering
    \caption{Implementación de SJF — reordenamiento por tamaño}
    \label{code:sjf-impl}
    \begin{lstlisting}[language=Java]
// SJFScheduler.java
@Override
public List<Path> ordenar(List<Path> pdfs) throws Exception {
    List<Path> ordenados = new ArrayList<>(pdfs);

    ordenados.sort(Comparator.comparingLong(pdf -> {
        try {
            return Files.size(pdf); // menor tamano = menor burst time estimado
        } catch (Exception e) {
            return Long.MAX_VALUE;  // al final si no se puede leer
        }
    }));

    return ordenados;
}
    \end{lstlisting}
\end{figure}

Si dos archivos tienen el mismo tamaño, el orden de llegada original actúa como
desempate, preservando la equidad entre procesos equivalentes. Una vez reordenada
la cola, la ejecución es idéntica a FCFS: cada archivo ocupa la CPU sin
interrupciones hasta completarse.

% ─────────────────────────────────────────────────────────────────
\subsubsection{Round Robin (RR)}

Round Robin es un algoritmo expropiativo que asigna a cada proceso un quantum de
tiempo fijo —en este sistema, 150\,ms. Si un archivo no termina dentro del quantum,
el hilo que lo procesa es interrumpido, el archivo se reencola al final de la cola
circular, y se retoma en la siguiente vuelta. Este ciclo continúa hasta que todos
los archivos completan su procesamiento.

La implementación utiliza un \texttt{Future} con \texttt{timeout} para simular la
expropiación:

\begin{figure}[H]
    \centering
    \caption{Implementación de Round Robin — quantum de tiempo con Future}
    \label{code:rr-impl}
    \begin{lstlisting}[language=Java]
// RoundRobinScheduler.java
private static final long QUANTUM_MS = 150;

Future<?> future = executor.submit(() -> {
    ProcesadorPDFTask tarea = new ProcesadorPDFTask(pdf);
    tarea.call();
});

try {
    future.get(QUANTUM_MS, TimeUnit.MILLISECONDS);
    // Termino dentro del quantum: registrar fin
    metrica.registrarFin();

} catch (TimeoutException e) {
    // Excedio el quantum: cancelar y reencolar
    future.cancel(true);
    cola.offer(pdf); // regresa al final de la cola circular
}
    \end{lstlisting}
\end{figure}

La cola circular se implementa con un \texttt{LinkedList} operado como
\texttt{Queue}. En cada vuelta se procesan todos los archivos pendientes;
los que exceden el quantum son reinsertados al final con \texttt{cola.offer()}.
El proceso se repite hasta que la cola queda vacía. Para evitar bucles infinitos
en caso de fallos persistentes, se establece un límite de 100 vueltas.
